% ReXGrounding Challenge paper: team method description (Appendix A)
% ---------------------------------------------------------------
% Please keep this to about half a page of text (300-450 words) plus one optional
% figure. Keep the headings below in this order; write "none" where something
% does not apply. We will typeset and may
% lightly edit for length and style. Citations: please list the reference at
% the end of the file in plain text; we will add them to the bibliography.

\subsection{<Team name> (rank <N>)}
\label{app:<teamname>}

% Everything the model saw in training, with sizes and licences. Include the
% challenge data, public datasets, in-house data, pretrained weights, and any
% derived labels (pseudo-labels, LLM-extracted text, anatomy masks). State
% explicitly whether any test-set scans, in image or text form, were used at any
% stage, including pretraining of components you did not train yourselves.
\paragraph{Data} <Challenge training set (2,992 scans) and validation set (200 scans); external data: ...; in-house data: ...; pretrained weights: ...; test-set exposure: none / describe.>

% Resampling, orientation, intensity windowing or normalisation, cropping, patch size.
\paragraph{Pre-processing} <...>

% Image backbone, text encoder, how the text conditions the image model, number
% of parameters if known. Name the base framework (for example VoxTell, nnU-Net).
\paragraph{Architecture and text encoder} <...>

% Loss, optimiser, learning-rate schedule, batch size, epochs or steps, data
% augmentation (image and text), hardware and training time, whether the
% validation set was used for training and how.
\paragraph{Training} <...>

% Number of models in the final submission, how they differ, how predictions were combined.
\paragraph{Ensembling} <...>

% Thresholding, component selection, anatomical restriction (lobe or side),
% any use of an LLM at inference.
\paragraph{Post-processing} <...>

% Average time per scan per finding at inference, and the hardware used.
\paragraph{Inference time and hardware} <...>

% Which submission is the final one (model name as shown on the leaderboard),
% and one or two sentences on what helped most and what did not help.
\paragraph{Final submission and lessons} <...>

% Link(s) to code and weights, with licence; or state that they are not released.
\paragraph{Code and weights availability} <...>

% Optional: one figure of the architecture or pipeline (vector PDF preferred).
% \begin{figure}[h]\centering\includegraphics[width=\textwidth]{<file>.pdf}\caption{<...>}\end{figure}

% References (plain text, we will format them):
% 1. ...
